\chapter{Evidence synthesis and meta-analysis}
\label{sec:p15-evidence-synthesis}

A combined estimate answers a question only when the constituent estimates
have an identified relation to that question. Precision cannot repair a
change of outcome, population, intervention, time horizon or comparison.
The evidence record of \ref{def:p19-evidence-record} supplies the study-level
information needed before numerical synthesis. The inferential distinction
between a confidence procedure and a posterior distribution remains the
one in \ref{def:p19-uncertainty-terms}.

The methodological presentation used here is Chapter~10 of the Cochrane
Handbook, version~6.5, last updated November~2024, Sections~10.1--10.14
\cite{DeeksEtAl2024Cochrane10}. The mathematical results below supply explicit
derivations and delimit what is exact when variances are known, what remains
a plug-in approximation when they are estimated, and what must be reported
when studies differ. The worked studies are constructed examples, not a
reanalysis of a clinical or psychological dataset.

\section{Specify the combined question}
\label{sec:p15-synthesis-compatibility}

\begin{hbdefinition}{def:p15-synthesis-record}{A synthesis record and a common target}
For each of \(k\geq1\) studies, record its population and recruitment,
intervention and comparator versions, outcome definition and time horizon,
assignment and analysis design, effect scale, estimate \(Y_i\), uncertainty
calculation, missingness, and participants or records shared with other
estimates. An \emph{estimand} \(\theta_i\) is the particular population
quantity that \(Y_i\) estimates under its stated design/model assumptions.
It is distinct from the estimate and from a causal interpretation.

A common-effect synthesis asserts \(\theta_i=\theta\) for every included
study after justified transformations. A finite weighted synthesis instead
targets the explicit quantity \(\theta_a=\sum_i a_i\theta_i\), where
the deterministic weights satisfy \(\sum_i a_i=1\). A random-effects
synthesis posits a distribution of underlying study effects, with mean
\(\mu\) and variance \(\tau^2\). These are three different targets.
Study inclusion does not by itself identify a target population distribution
for future studies or participants.
\end{hbdefinition}

\begin{hbmethod}{def:p15-synthesis-compatibility}{Compatibility before calculation}
Specify the review question and eligibility criteria, identify the actual
studies and overlapping reports, and check the data against their sources.
Align intervention/comparator, outcome, follow-up, effect direction, target
population and analysis design. Determine whether estimates are marginal or
conditional on different adjustment variables. Specify a common estimand or
the population distribution over which heterogeneous effects are to be averaged.
Declare the transformation and its assumptions for each scale. Retain the
study's clustering, repeated-measurement or crossover uncertainty calculation.
Then choose a pooling model. If no meaningful shared quantity is established,
present the separate estimates and their evidence records.
\end{hbmethod}

Clinical or task diversity can be relevant to a deliberately broad question;
it does not automatically forbid synthesis. Conversely, similar numerical
ranges do not establish a common measurement. In
\ref{def:p19-composite-comparability}, a slope for a latent log-odds score and
a slope for a proportion score have different units. Pooling them numerically
would require a justified transformation of both estimates and uncertainties.
Risk-of-bias differences can also produce heterogeneous observed estimates
without identifying genuine population effect modification
\cite[secs.~10.1, 10.10.1]{DeeksEtAl2024Cochrane10}.

\begin{hbdefinition}{def:p15-synthesis-effect-scales}{Effect scales and denominators}
For binary outcomes let \(p_T,p_C\) denote event probabilities for the
specified treatment and comparator, at the same defined horizon. The risk
difference is \(p_T-p_C\), the risk ratio is \(p_T/p_C\) when
\(p_C>0\), and the odds ratio is
\(\{p_T/(1-p_T)\}/\{p_C/(1-p_C)\}\) when both probabilities lie
strictly between zero and one. Positive finite ratios are conventionally pooled on their
natural-log scale and transformed back afterwards.

For quantitative outcomes, the mean difference compares means in the same
units. A standardized mean difference divides by a declared outcome standard
deviation; that denominator is part of the estimand, not merely a change of
display. A log mean ratio requires positive means and its own sampling model.
A proportional-odds effect for an ordinal scale requires a specified threshold
model. A rate ratio compares events per exposure time, while a hazard ratio
compares instantaneous hazards among those still at risk. Neither generally
equals a cumulative risk ratio.
\end{hbdefinition}

\begin{hbtheorem}{thm:p15-synthesis-scale}{proved=annex:p15.synthesis.scale}{Transformations that preserve or change the target}
The following identities make scale compatibility testable.
\begin{enumerate}
\item Applying the same affine transformation \(X'=b+cX\), with
      \(c>0\), to both groups multiplies a mean difference and its
      standard error by \(c\); a difference standardized by the correspondingly
      transformed outcome SD is unchanged. Changing the population SD
      independently of measurement units need not preserve that standardized target.
\item Replacing an event by its complement negates the risk difference and
      reciprocates the odds ratio. The complementary risk ratio is
      \((1-p_T)/(1-p_C)\), generally not the reciprocal risk ratio.
\item A common risk ratio \(r\) implies a risk difference
      \((r-1)p_C\). A common odds ratio \(o>0\) implies
      \(p_T=op_C/(1-p_C+op_C)\). Absolute effects therefore depend on
      comparator risk even when the relative effect is constant.
\item A pooled log ratio \(\sum_i a_i\log r_i\) back-transforms to
      \(\prod_i r_i^{a_i}\), not to the arithmetic mean of the ratios.
\end{enumerate}
\end{hbtheorem}

An odds ratio conditional on covariates is not generally the same estimand as
the marginal odds ratio. With equally frequent strata having comparator risks
\(1/10\) and \(1/2\), a conditional odds ratio of two in both strata
gives treatment risks \(2/11\) and \(2/3\). The marginal comparator
risk is \(3/10\) and the marginal treatment risk is \(14/33\); the
marginal odds ratio is \(98/57\), not two. This is a constructed
non-collapsibility example, not evidence of confounding in a particular study.
Its arithmetic is included in \ref{annex:p15.synthesis.scale}.

\section{Weighted estimates and exact uncertainty}
\label{sec:p15-synthesis-weights}

\begin{hbdefinition}{def:p15-synthesis-linear-model}{Fixed weights, covariance and sampling assumptions}
Let \(Y=(Y_1,\ldots,Y_k)^{\mathsf T}\) have finite second moments,
\(E(Y_i)=\theta_i\), and covariance matrix \(\Sigma\). A fixed
weight vector \(a\) is deterministic, chosen before the realized outcomes, and
satisfies \(\mathbf1^{\mathsf T}a=1\). Its linear estimate is
\(\widehat\theta_a=a^{\mathsf T}Y\). For independent estimates,
write \(v_i=\operatorname{Var}(Y_i)>0\),
\(w_i=1/v_i\), \(W=\sum_i w_i\), and
\(\widehat\theta=\sum_iw_iY_i/W\). These \(v_i\) are actual,
known sampling variances in the exact results. Reported squared standard
errors \(s_i^2\) are estimates of them and are different objects.
Throughout the interval calculations below, \(0<\alpha<1\).
\end{hbdefinition}

\begin{hbtheorem}{thm:p15-synthesis-linear}{proved=annex:p15.synthesis.linear}{Linear pooling, covariance and optimal common-effect weights}
For fixed \(a\),
\[
 E\widehat\theta_a=\sum_i a_i\theta_i,
 \qquad \operatorname{Var}(\widehat\theta_a)=a^{\mathsf T}\Sigma a.
\]
If \(\theta_i=\theta\) and \(\Sigma\) is positive definite, the
unique minimum-variance linear unbiased weight vector and its variance are
\[
 a_*={\Sigma^{-1}\mathbf1\over\mathbf1^{\mathsf T}\Sigma^{-1}\mathbf1},
 \qquad V_*={1\over\mathbf1^{\mathsf T}\Sigma^{-1}\mathbf1}.
\]
For independent estimates these reduce to \(a_i=w_i/W\) and
\(V_*=1/W\). If \(Y\) is multivariate normal and the covariance
is known, then \(a_*^{\mathsf T}Y\pm z_{1-\alpha/2}\sqrt{V_*}\)
is an exact confidence interval; in the independent case this is
\(\widehat\theta\pm z_{1-\alpha/2}/\sqrt W\). Its coverage is exactly
\(1-\alpha\) for \(\theta\), for
\(0<\alpha<1\), where \(z_p\) is a standard-normal quantile.
Without normality the same variance calculation is exact, but that normal
interval does not automatically have the stated coverage.
\end{hbtheorem}

The precision claim for inverse-variance weighting refers to this model and
class of linear estimators. Under heterogeneous fixed \(\theta_i\), the
same known-weight calculation estimates their precision-weighted average.
Its confidence interval concerns that particular average under the specified
sampling model; it does not supply a common effect or a prediction for another
population \cite[secs.~10.3.1, 10.10.4.1]{DeeksEtAl2024Cochrane10}.

\begin{hbtheorem}{thm:p15-synthesis-estimated-weights}{proved=annex:p15.synthesis.estimated-weights}{Random weights do not inherit a fixed-weight guarantee}
For integrable random normalized weights \(A_i(Y)\) and common mean
\(E(Y_i)=\theta\), whenever the products are integrable,
\[
 E\!\left[\sum_i A_iY_i\right]-\theta
       =\sum_i\operatorname{Cov}(A_i,Y_i).
\]
An estimator chosen by giving almost all weight to the larger of two
independent \(N(\theta,1)\) observations can therefore be biased despite
positive weights. In the limiting selection rule, its mean is
\(\theta+1/\sqrt\pi\). For a positive definite fixed covariance
matrix, misspecifying its entries also invalidates the nominal variance
unless the resulting weight vector and actual covariance are used together.
\end{hbtheorem}

Using \(1/s_i^2\) may give a useful large-sample procedure when study
estimates and standard errors are sufficiently well behaved. The exact theorem
above does not establish its bias or coverage: the weights can depend on the
same observations as the effects, and estimated small variances can be unstable.
Zero or undefined variance estimates must be investigated; they are not infinite
amounts of independent evidence. The approximate log-risk-ratio, log-odds-ratio
and standardized-difference calculations need their actual within-study
models, sample sizes and uncertainty formulas.

\section{Outcome-specific analysis and sparse evidence}
\label{sec:p15-synthesis-outcomes}

\begin{hbmethod}{def:p15-synthesis-outcome-methods}{Choose a procedure for the recorded outcome}
For a quantitative outcome on a common scale, pool comparable mean differences
with their design-correct standard errors. For different instruments measuring
the same intended construct, assess whether a standardized difference has a
common meaning; standardization cannot establish construct validity.

For binary outcomes distinguish inverse-variance, Mantel--Haenszel and Peto
procedures. Mantel--Haenszel pools arm-level tables with an effect-specific
weighting scheme. Peto's one-step method pools score/information quantities to
approximate a log odds ratio; it is not a generic risk-ratio procedure.
Sparse events, strongly unequal arms and large effects require particular
care with those approximations. Applying a fixed continuity correction to
zero cells changes estimates and their uncertainties, rather than recovering
unobserved events. A study with no events in either arm does not establish a
relative effect of one. Its absolute rarity information and review membership
must still be retained.

For ordinal data, use a common declared dichotomy, a justified quantitative
scale, or an appropriate proportional-odds model. For recurrent events retain
exposure time and within-person dependence. For time-to-event data retain the
actual censoring/analysis assumptions and distinguish log-rank from Cox
model summaries. A reported \(O-E\) score and its information produce a
one-step effect on the scale of that score's model. The notation alone does
not identify an odds ratio or hazard ratio.
\end{hbmethod}

These procedure distinctions are those of Sections~10.4--10.9
\cite{DeeksEtAl2024Cochrane10}; they do not assert that every procedure is
appropriate for every effect scale. A standard error must correspond to the
same estimate, rather than to a different adjusted model or outcome.

\begin{hbdefinition}{def:p15-synthesis-binary-tables}{Binary tables and the conditional odds model}
In study \(i\), let \(a_i,b_i,c_i,d_i\) be event/non-event counts
in the treatment and comparator arms, respectively. Put
\(n_{Ti}=a_i+b_i\), \(n_{Ci}=c_i+d_i\),
\(n_i=n_{Ti}+n_{Ci}\), and \(m_i=a_i+c_i\); both arms are nonempty.
The Mantel--Haenszel odds-ratio and risk-ratio estimates are
\[
 {\sum_i a_id_i/n_i\over\sum_i b_ic_i/n_i},\qquad
 {\sum_i a_in_{Ci}/n_i\over\sum_i c_in_{Ti}/n_i},
\]
when their denominators are positive. These are ratios of table-weighted
totals, not inverse-variance averages of study log ratios. Their uncertainty
is not supplied by the known-weight theorem for a different estimator.

For the conditional odds model, assume independent Bernoulli observations
within arms and independent studies, with a common log odds ratio \(\delta\)
but separate comparator risks strictly between zero and one. Conditional on the arm sizes and total events,
the treatment event count \(A_i\) has mass proportional to
\[
 {n_{Ti}\choose x}{n_{Ci}\choose m_i-x}e^{\delta x},
 \quad \max(0,m_i-n_{Ci})\leq x\leq\min(n_{Ti},m_i).
\]
Write \(Z_i(\delta)\) for the finite sum of these weights. At
\(\delta=0\), define
\[
 E_i={n_{Ti}m_i\over n_i},\qquad
 V_i={n_{Ti}n_{Ci}m_i(n_i-m_i)\over n_i^2(n_i-1)}.
\]
The table formulas and Peto calculation are recorded in the May~2022
formula supplement \cite[pp.~2, 4--5, 7]{DeeksHiggins2022RevMan}; this older
formula source is used here for those specific binary calculations. The
November~2024 chapter supplies the distinct current random-effects discussion.
Peto's binary one-step calculation is
\(\widehat\delta_P=\sum_i(a_i-E_i)/\sum_iV_i\), with approximate
SE \((\sum_iV_i)^{-1/2}\), when \(\sum_iV_i>0\).
\end{hbdefinition}

\begin{hbtheorem}{thm:p15-synthesis-score}{proved=annex:p15.synthesis.score}{The score calculation and its approximation boundary}
The conditional log likelihood has score
\(U(\delta)=\sum_i\{a_i-E_\delta(A_i)\}\) and negative second
derivative \(I(\delta)=\sum_i\operatorname{Var}_\delta(A_i)\).
At zero, its score/information are \(\sum_i(a_i-E_i)\) and
\(\sum_iV_i\). The displayed Peto estimate is one Newton step from
zero towards a conditional maximum-likelihood solution. Under the conditional
null its score has exactly mean zero and variance \(I(0)\); a normal
confidence interval for the one-step estimate remains an approximation.
Studies with no events in either arm have zero score and zero information
in this conditional relative-effect model. Their exclusion from that
calculation does not establish zero risk or no absolute treatment effect.
\end{hbtheorem}

\begin{hbtheorem}{thm:p15-synthesis-table-delta}{proved=annex:p15.synthesis.table-delta}{A declared large-arm uncertainty calculation for table ratios}
For independent Bernoulli arms, suppose the number of studies is fixed,
\(n_{gi}/n\to\lambda_{gi}>0\) as \(n\to\infty\), and every
event probability is fixed strictly between zero and one. If the odds ratio
is common, the Mantel--Haenszel odds-ratio estimate is consistent for it.
If the risk ratio is common, the corresponding risk-ratio estimate is
consistent for it. A delta-method variance calculation for either log ratio
\(f=\log(R/S)\) is
\[
 \widehat V_f=\sum_i\left\{
       f_{a_i}^2 n_{Ti}\widehat p_{Ti}(1-\widehat p_{Ti})
      +f_{c_i}^2 n_{Ci}\widehat p_{Ci}(1-\widehat p_{Ci})\right\},
\]
where \(\widehat p_{gi}\), \(g\in\{T,C\}\), is the observed arm
event proportion: \(\widehat p_{Ti}=a_i/n_{Ti}\) and
\(\widehat p_{Ci}=c_i/n_{Ci}\). The derivatives are evaluated
at the observed counts. For the odds-ratio totals
\(R=\sum a_i(n_{Ci}-c_i)/n_i\),
\(S=\sum(n_{Ti}-a_i)c_i/n_i\), they are
\[
 f_{a_i}={n_{Ci}-c_i\over n_iR}+{c_i\over n_iS},\qquad
 f_{c_i}=-{a_i\over n_iR}-{n_{Ti}-a_i\over n_iS}.
\]
For the risk-ratio totals \(R=\sum a_in_{Ci}/n_i\),
\(S=\sum c_in_{Ti}/n_i\), they are
\(f_{a_i}=n_{Ci}/(n_iR)\), \(f_{c_i}=-n_{Ti}/(n_iS)\).
The standardized log-estimation error is asymptotically standard normal
using \(\widehat V_f\). This is a specified independent-arm, large-arm
delta procedure, not an assertion of exact small-sample coverage or identity
with every finite-sample variance formula in the software supplement.
\end{hbtheorem}

\begin{hbtheorem}{thm:p15-synthesis-change}{proved=annex:p15.synthesis.change}{Changes, final values and a declared conversion model}
Let \(X_0,X_1\) be baseline and final measurements. Then
\[
 \operatorname{Var}(X_1-X_0)
 =\operatorname{Var}(X_1)+\operatorname{Var}(X_0)
                   -2\operatorname{Cov}(X_0,X_1).
\]
If randomized groups have equal population baseline means, their difference
in mean changes equals their difference in final means. Their estimated
precision need not be equal. Standardizing by a change SD and by a final-value
SD generally gives different targets.

For a continuous latent outcome with equal-scale logistic laws
\(F_g(x)=\{1+\exp(-(x-m_g)/s)\}^{-1}\), \(s>0\), define the
event as exceeding a common threshold. Its log odds ratio is
\((m_T-m_C)/s\). If that logistic scale has outcome SD
\(s\pi/\sqrt3\), the standardized mean difference is
\((\sqrt3/\pi)\log\mathrm{OR}\). The conversion requires this
latent distribution and equal-variability model; a shared outcome label alone
does not establish it. A standard error transforms by the same constant only
within the assumed conversion.
\end{hbtheorem}

Post-intervention and change-score MDs may thus estimate the same randomized
comparison while their SDs remain different. The required standard errors
use the corresponding data. Post-intervention and change-score SMDs cannot
be combined merely because both are called standardized differences. Medians,
means, arithmetic means of raw values and means of logarithms are also distinct.
Skewed small-sample data can make mean/SD summaries inadequate. A bounded-scale
mean-to-SD diagnostic is a warning signal, not proof of normality; transforming
or retrieving appropriate summaries may be necessary
\cite[secs.~10.5--10.6]{DeeksEtAl2024Cochrane10}.

\section{Heterogeneous effects and the uncertainty in heterogeneity}
\label{sec:p15-synthesis-heterogeneity}

\begin{hbdefinition}{def:p15-synthesis-random-model}{The normal random-effects model}
For \(k\geq2\), posit
\[
 \Theta_i\sim N(\mu,\tau^2),\qquad
 Y_i=\Theta_i+e_i,\qquad e_i\sim N(0,v_i),
\]
with all effects and errors independent, \(v_i>0\), and
\(\tau^2\geq0\). The degenerate \(\tau^2=0\) case has
\(\Theta_i=\mu\). For a specified heterogeneity variance \(u\geq0\),
write
\[
 w_i(u)={1\over v_i+u},\quad W(u)=\sum_iw_i(u),\quad
 M(u)={\sum_iw_i(u)Y_i\over W(u)},\quad
 Q(u)=\sum_iw_i(u)(Y_i-M(u))^2.
\]
The source's random-effects summary concerns the mean of a modelled
distribution of study effects. Systematic bias is not made mean zero by naming
it a random effect. Generalization requires a reason to relate that modelled
distribution to the intended setting.
\end{hbdefinition}

\begin{hbtheorem}{thm:p15-synthesis-random-known}{proved=annex:p15.synthesis.random-known}{Known heterogeneity, mean estimation and prediction}
Under the preceding model with known \(v_i\) and known \(\tau^2\),
\(M(\tau^2)\) is unbiased for \(\mu\), normally distributed with
variance \(1/W(\tau^2)\). Its normal confidence interval for the mean
has exact coverage. For an independent new effect
\(\Theta_*\sim N(\mu,\tau^2)\),
\[
 M(\tau^2)\ \pm\ z_{1-\alpha/2}
            \sqrt{\tau^2+1/W(\tau^2)}
\]
is an exact \(1-\alpha\) prediction procedure for \(\Theta_*\).
For a future noisy estimate with independent normal error of variance \(v_*\),
add \(v_*\) under the square root. A confidence interval for \(\mu\)
alone is not this prediction interval.
\end{hbtheorem}

For equal within-study variances \(v_i=v\), the variance of the estimated
mean is \((v+\tau^2)/k\), tending to zero as the number of studies grows.
The future-effect prediction variance is \(\tau^2+(v+\tau^2)/k\),
tending to \(\tau^2\). Precision about the mean can therefore increase
while genuine effect variation remains.

\begin{hbtheorem}{thm:p15-synthesis-q}{proved=annex:p15.synthesis.q}{Dispersion, moment estimation and a known-variance confidence set}
Put \(w_i=w_i(0)\), \(W=W(0)\), and
\(C=W-\sum_iw_i^2/W>0\). Under independent effects/errors with
the random model's means and variances, without requiring normality,
\[
                    E Q(0)=k-1+\tau^2C.
\]
The moment estimate is therefore
\[
 \widehat\tau_{\rm DL}^2
    =\max\{0,(Q(0)-(k-1))/C\}.
\]
The truncated estimate is not asserted to be unbiased. Under the full normal
model and known \(v_i\), \(Q(\tau^2)\sim\chi^2_{k-1}\), and
\(Q(u)\) is nonincreasing for \(u\geq0\). Hence
\[
 \{u\geq0:q_{\alpha/2,k-1}\leq Q(u)
                         \leq q_{1-\alpha/2,k-1}\}
\]
has exact \(1-\alpha\) coverage for \(\tau^2\); here
\(q_{p,d}\) is a chi-squared quantile. The set may be empty for a
realized dataset. Replacing the \(v_i\) by estimated variances does not
preserve this exact coverage claim automatically.
\end{hbtheorem}

\begin{hbtheorem}{thm:p15-synthesis-reml}{proved=annex:p15.synthesis.reml}{The restricted normal likelihood}
In the known-within-variance normal model, a restricted log likelihood for
\(u=\tau^2\geq0\), up to an additive constant, is
\[
 \ell_R(u)=-\tfrac12\left\{
       \sum_i\log(v_i+u)+\log W(u)+Q(u)\right\}.
\]
Its derivative is
\[
 \ell_R'(u)=\tfrac12\left\{
  \sum_iw_i(u)^2(Y_i-M(u))^2
       -\sum_iw_i(u)+{\sum_iw_i(u)^2\over W(u)}\right\}.
\]
A global maximizer on \([0,\infty)\) exists when \(k\geq2\).
It may be at zero; uniqueness is not asserted. This is not the profile
likelihood obtained simply by substituting the best-fitting mean into the
full likelihood. Practical REML with estimated within-study variances remains
a plug-in normal-model procedure.
\end{hbtheorem}

The November~2024 presentation distinguishes REML from the
DerSimonian--Laird moment procedure and treats uncertainty in the estimated
between-study variance explicitly. An estimated \(\tau^2=0\), a
nonsignificant dispersion test, or a small \(I^2\) does not establish that
all effects are equal. Few studies give little information about a variance.
A dispersion test with many precise studies can detect differences too small
to matter to the review question
\cite[secs.~10.10.2, 10.10.4.4--10.10.4.5]{DeeksEtAl2024Cochrane10}.

\begin{hbdefinition}{def:p15-synthesis-inconsistency}{Two related inconsistency summaries}
The Q-based summary is
\[
 I_Q^2=100\max\{0,(Q(0)-(k-1))/Q(0)\}\ \%,
\]
with zero assigned when \(Q(0)=0\). Put
\(\widetilde v=(k-1)/C\). The tau-based summary is
\[
 I_\tau^2=100{\widehat\tau^2\over\widehat\tau^2+\widetilde v}\ \%.
\]
Here the factor 100 gives the percentage number: for example, a ratio of
\(1/2\) is 50 percent. With the DL estimate the two summaries are identical.
Write \(I_{\rm DL}^2\) for the tau-based summary using DL.
With another heterogeneity estimator, such as REML, they need not be identical.
They describe dispersion relative to within-study precision, not the proportion
of studies that are biased or the probability that heterogeneity is present.
\end{hbdefinition}

\begin{hblemma}{thm:p15-synthesis-inconsistency}{proved=annex:p15.synthesis.inconsistency}{The moment-based identity and a precision limit}
For the definitions above, \(I_{\rm DL}^2=I_Q^2\). For a fixed
positive \(\widehat\tau^2\), the tau-based percentage tends to 100
as \(\widetilde v\downarrow0\), without any increase in the absolute
heterogeneity variance. Thus a universal magnitude interpretation from
\(I^2\) alone is unavailable.
\end{hblemma}

\section{Estimated intervals and posterior uncertainty}
\label{sec:p15-synthesis-intervals}

\begin{hbdefinition}{def:p15-synthesis-hksj}{Wald and HKSJ calculations}
After estimating \(u\) and replacing \(v_i\) by the reported
\(s_i^2\) where needed, calculate \(w_i(u),W(u),M(u)\) by the
same formulas. The Wald calculation uses \(s_M^2=1/W(u)\) and
\(M(u)\pm z_{1-\alpha/2}s_M\). The HKSJ calculation uses
\[
 q_H={\sum_iw_i(u)(Y_i-M(u))^2\over k-1},\qquad
 s_H^2={q_H\over W(u)},\qquad
 M(u)\pm t_{k-1,1-\alpha/2}s_H.
\]
A variance safeguard sometimes replaces \(q_H\) by \(\max\{1,q_H\}\);
that is a separate modified procedure. Estimated weights and estimated
heterogeneity make these model-dependent approximate confidence calculations,
not instances of the exact known-variance theorem. In particular HKSJ can
be narrower than Wald when \(q_H\) is small, and its t multiplier can be
very large with two or three studies.
\end{hbdefinition}

\begin{hbtheorem}{thm:p15-synthesis-fixed-shape-t}{proved=annex:p15.synthesis.fixed-shape-t}{When a weighted Student interval really is exact}
For \(k\geq2\), if deterministic positive weights \(w_i\) are known and the observations
are independent \(N(\mu,c/w_i)\) with a common unknown \(c>0\),
then \(q_H/c\sim\chi^2_{k-1}/(k-1)\), independently of the weighted
mean, and \((M-\mu)/\sqrt{q_H/W}\sim t_{k-1}\). The displayed
weighted Student interval is therefore exact under this fixed-shape model.
Choosing weights by estimating \(\tau^2\) from those same observations
does not satisfy that premise in general.
\end{hbtheorem}

The chapter's simple approximate prediction calculation is
\(M\pm t_{k-1,1-\alpha/2}\sqrt{\widehat\tau^2+s_M^2}\).
Its implementation discussion also distinguishes using the normal multiplier
when Wald is selected and the t multiplier when HKSJ is selected; the associated
summary standard error must match the chosen procedure. Neither version is
an exact prediction guarantee for estimated heterogeneity. Report the
estimator, interval method, number of studies, heterogeneity estimate and its
uncertainty together. More studies can make the mean precise while genuine
effect variation stays large \cite[secs.~10.10.4.2--10.10.4.5]{DeeksEtAl2024Cochrane10}.

\begin{hbtheorem}{thm:p15-synthesis-bayes}{proved=annex:p15.synthesis.bayes}{Proper-prior normal synthesis and predictive variance}
In the normal model, give \((\mu,u)\), \(u\geq0\), a proper
prior probability measure \(\Pi_{\rm syn}\) with finite second moment for \(\mu\) and finite
mean for \(u\). The posterior has kernel
\[
 \Pi_{\rm syn}(d\mu,du)\prod_i(v_i+u)^{-1/2}
       \exp\!\left[-\tfrac12\sum_i{(Y_i-\mu)^2\over v_i+u}\right].
\]
Its normalizing constant is finite and strictly positive. For an independent
new effect from the model,
\[
 E(\Theta_*\mid Y)=E(\mu\mid Y),\qquad
 \operatorname{Var}(\Theta_*\mid Y)
    =E(u\mid Y)+\operatorname{Var}(\mu\mid Y).
\]
The prior and model, including their heterogeneity uncertainty, are part of
this posterior statement. Reusing the same study evidence in the prior and
likelihood without modelling that dependence counts it twice.
\end{hbtheorem}

Bayesian intervals concern posterior mass, with sensitivity to priors and
likelihood assumptions; frequentist coverage concerns repetition under a
specified sampling model. A uniform density over an unbounded parameter
range is not a proper probability distribution. An improper-prior analysis
requires its own posterior-propriety argument. The mathematical proposition
here uses a proper prior and supplies that argument
\cite[sec.~10.13]{DeeksEtAl2024Cochrane10}.

\section{Dependence, subgroups and missing information}
\label{sec:p15-synthesis-dependence}

\begin{hbtheorem}{thm:p15-synthesis-dependence}{proved=annex:p15.synthesis.dependence}{Duplicated records and a shared comparator}
Two reports of the same unbiased estimate \(Y\) with variance \(v>0\)
have average \(Y\), with variance \(v\), not \(v/2\). If independent
treatment-arm sample means share one comparator mean, their contrasts have
covariance equal to the comparator mean's variance. For two common-mean
estimates of variance \(v\) and correlation \(r\in[-1,1]\), the
equally weighted mean has variance \(v(1+r)/2\). Assuming independence
without evidence can thus substantially understate or overstate uncertainty.
\end{hbtheorem}

Use the study/participant as the appropriate independent unit. Link multiple
reports from the same study, choose a predeclared outcome/time-point rule,
combine compatible arms within the original design, or model their covariance.
Splitting a shared comparator's sample size among contrasts may be a particular
approximation; it does not make the original observations independent.
An unknown covariance needs a justified model or a sensitivity range constrained
to valid covariance matrices. Analyses of cluster-randomized and crossover
studies must retain the within-study dependence adjustment before their estimates
are combined. One cannot infer a covariance matrix from separate standard errors
alone \cite[secs.~10.3.3, 10.12.1]{DeeksEtAl2024Cochrane10}.

\begin{hbtheorem}{thm:p15-synthesis-subgroups}{proved=annex:p15.synthesis.subgroups}{A subgroup contrast is a separate comparison}
Partition independent known-variance estimates into \(G\geq2\) nonempty
groups. Write \(W_g=\sum_{i\in g}w_i\),
\(M_g=\sum_{i\in g}w_iY_i/W_g\). Then
\[
 Q(0)=\sum_g\sum_{i\in g}w_i(Y_i-M_g)^2
                    +\sum_gW_g(M_g-M(0))^2.
\]
Under a common normal effect, the between-group term has
\(\chi^2_{G-1}\) law. For two group summaries with actual variances
\(V_1,V_2\) and covariance \(c_{12}\), the difference has variance
\(V_1+V_2-2c_{12}\). A significant result within one group and a
nonsignificant result within another does not establish a group difference.
\end{hbtheorem}

For example, constructed independent group summaries \(1.0\) with SE
\(0.4\) and \(0.9\) with SE \(0.5\) have null standardized
statistics 2.5 and 1.8, on opposite sides of the usual 1.96 threshold.
Their difference is 0.1 with SE \(\sqrt{0.41}\), giving only about
0.156 standardized units. Comparing the two significance labels would
claim an interaction unsupported by the actual contrast.

Study-level subgroup analyses are observational comparisons even when each
trial randomized its own participants. Meta-regression needs a predeclared
covariate model, enough studies and covariate variation, appropriate weights,
and residual heterogeneity. Confounded covariates can make explanations
non-identifiable. Equal study mean ages do not rule out within-study age
modification: a model with effects increasing with age and a model with a
constant effect can have the same study-level averages. Prognostic factors
need not modify an intervention effect. Dose-response conclusions from
randomized within-study dose contrasts have a different evidential standing
from associations across study doses. Searching many covariates or cutpoints
changes the false-positive opportunity; a selected explanation requires
confirmation \cite[sec.~10.11]{DeeksEtAl2024Cochrane10}.

\begin{hbdefinition}{def:p15-synthesis-missing}{Missingness and the retained denominator}
Missing studies, unreported outcomes, missing summary variances, missing
participant outcomes and missing moderator values are different problems.
For participant data, let \(R\) record what is observed. Missing completely
at random means independence of \(R\) from the complete data; missing at
random means that, conditional on observed data and covariates, the missingness
law does not depend on the unobserved values. Neither statement by itself
justifies every complete-case analysis. The required model and conditioning
variables must be supplied. Missing not at random permits dependence on the
unobserved values after that conditioning.
\end{hbdefinition}

\begin{hbtheorem}{thm:p15-synthesis-missing-bounds}{proved=annex:p15.synthesis.missing-bounds}{Finite missing-outcome bounds}
In a group with \(N=n+m>0\) allocated participants, \(n\) observed
outcomes, \(m\) missing outcomes and \(e\) observed events,
the full-sample event proportion lies in
\([e/N,(e+m)/N]\). For two groups its risk difference therefore lies in
\[
 \left[{e_T\over N_T}-{e_C+m_C\over N_C},
             {e_T+m_T\over N_T}-{e_C\over N_C}\right].
\]
Both endpoints are attainable by assigning the missing binary outcomes.
These are deterministic bounds for the allocated sample, not population
confidence limits. A sensitivity model assigning event risks to the missing
participants supplies additional assumptions, whose uncertainty must be retained.
\end{hbtheorem}

Request missing information where possible, retain eligible studies in the
review even when they cannot enter a numerical synthesis, and state every
imputation assumption. Filling a missing SD or event count with a chosen
number does not make it observed. Treating imputed quantities as known can
understate uncertainty. Single imputation, multiple imputation with uncertainty,
and a full missing-data model have different inferential requirements.
The separate assessment of selectively unavailable studies/results belongs
to missing-evidence methodology; a pooled estimate alone cannot detect or
correct that selection \cite[sec.~10.12]{DeeksEtAl2024Cochrane10}.

\section{A complete pooling calculation and sensitivity}
\label{sec:p15-synthesis-example}

\begin{hbtheorem}{thm:p15-synthesis-sensitivity}{proved=annex:p15.synthesis.sensitivity}{Fixed-weight sensitivity identities}
For positive deterministic \(w_i\), let \(W=\sum_iw_i\) and
\(M=\sum_iw_iY_i/W\). If \(W>w_j\), deleting study \(j\)
without changing the other weights gives
\[
 M_{-j}-M={w_j\over W-w_j}(M-Y_j),\qquad
 {\partial M\over\partial w_j}={Y_j-M\over W}.
\]
If study bias shifts satisfy \(|b_i|\leq B_i\), the pooled shift
is bounded by \(\sum_i(w_i/W)B_i\), with both extremes attainable.
Deleting a study and re-estimating heterogeneity changes all weights, so the
deletion identity alone does not compute that random-effects sensitivity.
\end{hbtheorem}

\begin{hbexample}{def:p15-synthesis-complete-example}{Three compatible studies and two stated models}
Suppose three independent estimates on the same difference scale are
\(Y=(-2,-1,2)\), with \emph{known} variances \(v=(1,1,4)\).
No actual study is claimed. The common-effect weights are
\((1,1,1/4)\), \(W=9/4\), and
\[
 M(0)=-10/9\approx-1.1111,\quad
 \operatorname{Var}M(0)=4/9,\quad Q(0)=29/9.
\]
Under the common normal model, the 95 percent mean interval is approximately
\([-2.4178,0.1955]\). The data also permit a heterogeneous interpretation;
the dispersion estimate does not select a model automatically.
Here \(C=4/3\) and \(\widehat\tau_{\rm DL}^2=11/12\).
The resulting random-effects plug-in weights are
\((12/23,12/23,12/59)\), with
\[
 W=1692/1357,\qquad M=-131/141\approx-0.9291,
 \qquad s_M^2=1357/1692\approx0.8020.
\]
The plug-in Wald interval is approximately \([-2.6843,0.8262]\).
It is not the exact known-heterogeneity interval because \(11/12\) was
estimated from these three observations. The HKSJ and prediction calculations,
as well as removal of the third estimate, are completed in
\ref{annex:p15.synthesis.example}; with only three studies their uncertainty
must be interpreted cautiously.

If the third estimate is actually another report from the same participants
as the first, these independence weights are invalid. If it measures a different
construct or time horizon without an identified bridge, the numerical mean
answers no established shared question. If its positive sign makes it an
unwanted outlier, removing it for that reason would select the conclusion;
report the declared sensitivity and its rationale.
\end{hbexample}

\begin{hbtheorem}{thm:p15-synthesis-example}{proved=annex:p15.synthesis.example}{Worked intervals and failure constructions}
The displayed three-study calculation follows from the stated formulas.
The same formulas produce distinct common-effect, estimated-heterogeneity,
HKSJ and prediction summaries. A duplicated estimate falsely treated as
independent understates its standard error by a factor \(\sqrt2\).
The displayed group estimates give incompatible significance labels without
a significant subgroup contrast. Twenty independent level-0.05 null comparisons
have probability \(1-0.95^{20}\approx0.642\) of at least one rejection;
without independence the union bound is at most one. A pair of study-level
average effects does not identify an unreported within-study effect modifier.
\end{hbtheorem}

\begin{hbmethod}{def:p15-synthesis-report}{Report the synthesis and the decisions that can change it}
State the target and compatible study set, effect scale and direction,
independent units and overlap handling, estimate/SE extraction, weighting
procedure, heterogeneity estimator, interval method and any prediction target.
Report the constituent effects with their uncertainties, the number of studies,
the pooled estimate and interval, absolute heterogeneity and its uncertainty,
and the extent of missing information. Keep confidence and prediction intervals
distinct. Where appropriate, translate the result to an absolute scale at a
stated comparator risk, retaining uncertainty about that risk.

Compare justified alternative eligibility boundaries, outcome times, bias
assessments, missing-data assumptions, variance/correlation imputations,
effect measures, common/random-effects procedures, heterogeneity estimators
and small-sample intervals. Specify prior choices in a Bayesian analysis.
A sensitivity analysis examines alternative ways of estimating the declared
target; a subgroup analysis estimates and compares distinct group targets.
Report important changes in magnitude, sign and uncertainty together with
the assumptions producing them. Agreement across the chosen analyses does
not establish robustness to an unexamined selection or measurement failure.
\end{hbmethod}

\begin{proofboundary}{def:p15-synthesis-boundary}{Limits of the combined conclusion}
The exact normal, fixed-weight and covariance results require their explicit
premises. Practical estimated weights, sparse-event approximations, unknown
heterogeneity, selective availability and unmeasured biases retain their own
uncertainties. A random-effects mean does not correct a systematically biased
corpus. A nonsignificant heterogeneity or interaction test does not establish
absence. The chapter supplies no causal transport theorem from an average
across observed studies to every new population, and no general guarantee
that a selected subgroup explanation is the cause of observed differences.
\end{proofboundary}
